%%%%%%%%%%%%%%%%%%%%%%%%%%%%%%%%%%%%%%%%%%%%%%%%%%%%
\section{Introduction}

Autonomous AI systems and robotic agents generate streams of sensor events
that encode behavioral patterns.
Understanding \textit{when} these patterns densely co-occur --- not merely whether
they exist globally --- is essential for real-time monitoring, anomaly detection,
and adaptive control in intelligent environments.
In multi-agent systems and sensor-equipped robotic platforms,
detecting that a specific set of events repeatedly clusters within particular
time intervals can reveal coordinated behaviors, routine structures,
or emerging anomalies that demand timely response.
The approach also has broad applicability across domains such as retail and epidemiology.

We propose \textit{Streaming Apriori-Window}, a dense interval mining
algorithm for event streams parameterized by a scan frequency $n$.
The algorithm identifies all time intervals in which itemset occurrence
frequency within a fixed-width sliding window meets a threshold.
Two speedup techniques --- stride-based skipping and Apriori-style
candidate pruning --- are incorporated to achieve practical processing speed.

The main contributions of this paper are as follows:
\begin{itemize}
  \item We formalize the streaming dense interval mining problem with
        an explicit scan frequency parameter $n$ that controls the tradeoff
        between output latency and processing cost.
  \item We prove a confirmed output guarantee: intervals ending before
        the current scan position are unaffected by future events.
  \item We evaluate the algorithm on synthetic and real-world CASAS
        smart home sensor data, demonstrating practical processing speed
        and interpretable outputs.
\end{itemize}
